\documentclass[12pt]{article}
\usepackage[T1]{fontenc}
\usepackage[utf8]{inputenc}
\usepackage{lmodern}
\usepackage{textcomp}
\usepackage{enumitem}
\usepackage{geometry}
\usepackage{graphicx}
\usepackage{amsmath, amssymb}
\geometry{margin=1in}
\begin{document}
\begin{center}
Chemistry - Undergraduate Level Assessment 16 \
Total Marks: 40 \
Answer all questions.
\end{center}

\section*{Multiple Choice (10 marks)}
\begin{enumerate}[label=\arabic*.]
\item What is the NMR frequency of 31 P in a 20.0 T magnetic field?
\begin{enumerate}[label=(\alph*)]
    \item 54.91 MHz
    \item 239.2 MHz
    \item 345.0 MHz
    \item 2167 MHz
\end{enumerate}
\item Which of the following statements about the lanthanide elements is NOT true?
\begin{enumerate}[label=(\alph*)]
    \item The most common oxidation state for the lanthanide elements is +3.
    \item Lanthanide complexes often have high coordination numbers (\textgreater{} 6).
    \item All of the lanthanide elements react with aqueous acid to liberate hydrogen.
    \item The atomic radii of the lanthanide elements increase across the period from La to Lu.
\end{enumerate}
\item Which of the following statements most accurately explains why the 1 H spectrum of 12 CHCl 3 is a singlet?
\begin{enumerate}[label=(\alph*)]
    \item Both 35 Cl and 37 Cl have I = 0.
    \item The hydrogen atom undergoes rapid intermolecular exchange.
    \item The molecule is not rigid.
    \item Both 35 Cl and 37 Cl have electric quadrupole moments.
\end{enumerate}
\item Which of the following is always true of a spontaneous process?
\begin{enumerate}[label=(\alph*)]
    \item The process is exothermic.
    \item The process does not involve any work.
    \item The entropy of the system increases.
    \item The total entropy of the system plus surroundings increases.
\end{enumerate}
\item Calculate the equilibrium polarization of 13 C nuclei in a 20.0 T magnetic field at 300 K.
\begin{enumerate}[label=(\alph*)]
    \item 10.8 x 10^-5
    \item 4.11 x 10^-5
    \item 3.43 x 10^-5
    \item 1.71 x 10^-5
\end{enumerate}
\end{enumerate}

\section*{True / False (10 marks)}
\textit{Answer True or False}
\begin{enumerate}[label=\arabic*.]
\setcounter{enumi}{5}
\item True or False: The thermal stability of group-14 hydrides increases in the order: CH 4 \textless{} SiH 4 \textless{} GeH 4 \textless{} SnH 4 \textless{} PbH 4.
\item True or False: The orbital angular momentum quantum number, l, of the electron most easily removed from ground-state aluminum is 2.
\item True or False: The +1 oxidation state is more stable than the +3 oxidation state for thallium (Tl).
\item True or False: The magnetogyric ratio of 23 Na is 7.081 x $10^7$ T^-1 s^-1.
\item True or False: The difference in infrared absorption frequency between DCl and HCl is primarily due to their reduced mass.
\end{enumerate}

\section*{Long Form Response (20 marks)}
\textit{Answer each question with a clear and complete explanation.}
\begin{enumerate}[label=\arabic*.]
\setcounter{enumi}{10}
\item Discuss the role of metal ions as paramagnetic quenchers in chemical reactions, specifically analyzing why Cr$_{3}$+ is not suitable for this purpose.
\item Discuss the solid-state structures of the principal allotropes of elemental boron, focusing on the role of B$_{12}$ icosahedra in their formation and stability.
\end{enumerate}

\vfill
\noindent\textit{End of Paper}
\end{document}